\section{A complete canonical phase-estimation comparator}
\label{sec:canonical-qpe-comparator}

I use the established estimator of \citet{brassard}, Section 4, Algorithm
\texttt{Est\_Amp} and Theorem 12, as an additional quantum comparator. The
source actually consulted is the AMS-copyright typeset 2002 chapter, hosted
by the volume editor; its amplitude-dependent bound is an established
result. The contribution of this section is its checked transfer to the
finite financial flags, prescribed confidence, outward output arithmetic
and complete emitted gate profiles. I do not claim a new phase-estimation
algorithm or infer physical advantage from a query count.

\subsection{Logical hypotheses and controlled Grover phase}

Let $A$ denote the exact logical flag preparation on the finite input and
flag register, with all private work initially zero. Let $P_1$ project onto
flag one and set $p=\langle0|A^\dagger P_1 A|0\rangle$. The proved prior is
$0\leq p\leq u\leq1$. Each selected financial flag must retain its own
finite encoding and deterministic price certificate. Write
$S_1=I-2P_1$, $S_0=I-2|0\rangle\langle0|$ on the active logical register,
and
\[
 Q=-A S_0 A^\dagger S_1.
\]
These operators are logical unitaries. The measured clean-AND implementation
is admissible only through the independently proved instrument statement:
for every complete outcome string $z$, the corrected, reset implementation
has Kraus operator $K_z=c_z A$ on the specified clean-work subspace,
$\sum_z|c_z|^2=1$, with $c_z$ independent of the logical state and of any
external reference. The same statement is required for the separately
implemented logical adjoint. Summing branches gives
\[
 \sum_z(K_z\otimes I)\rho(K_z^\dagger\otimes I)
 =(A\otimes I)\rho(A^\dagger\otimes I)
\]
for every joint logical-reference state $\rho$. Thus each instrument can be
replaced by its exact logical channel in the proof. The physical tape
contains measurements and is not described as a measurement-free unitary.
This transfer requires the full scalar-Kraus statement; a basis-state truth
table or flag probability alone would not suffice.

For phase control $c$, the controlled iterate is emitted chronologically as
\[
 \mathrm{CZ}_{c,f},\quad A^\dagger,\quad
 \mathrm{C}_c(S_0),\quad A,\quad Z_c.
\]
When $c=0$, the reflections are identities and the two unconditional
logical channels cancel. When $c=1$, their product is
$-A S_0 A^\dagger S_1=Q$. The last $Z_c$ is essential: a global minus on an
uncontrolled iterate becomes a relative phase between the two control
states. Removing it changes the estimated eigenphase by one half. The
coherent replay with $p=1$ gives a direct counterexample: the correct phase
outcome is $M/2$, whereas omitting the conditional sign gives zero.
Unconditional financial arithmetic is therefore reused inside controlled
$Q$; no arbitrary promotion of every arithmetic gate to an extra-controlled
gate is assumed.

For an active register of $n\geq2$ wires, the clean conjunction construction
for $S_0$ uses $2(n-2)$ exact Toffolis and one final CZ. Its controlled
version uses $2(n-1)$ exact Toffolis and one final CZ, with the phase-control
wire not negated by the surrounding $X$ gates. The controlled good
reflection replaces the original flag $Z$ by CZ, and $Z_c$ supplies the
conditional minus. Consequently the additional exact-Toffoli charge is
two Toffolis, or fourteen T gates under the specified seven-T expansion.
The actual emitted controlled profile, including its Clifford gates,
scratch and serial depth, is charged rather than treating this T increment
as a complete resource description. The corrected measurement instruments
inside $A$ and its adjoint retain their own observation, feed-forward and
reset charges.

\subsection{Phase law and amplitude error}

For $0<p<1$, let $|g\rangle=P_1A|0\rangle/\sqrt p$ and
$|b\rangle=(I-P_1)A|0\rangle/\sqrt{1-p}$. Define
$\theta=\arcsin\sqrt p\in(0,\pi/2)$. On the ordered basis
$(|g\rangle,|b\rangle)$,
\[
 Q=\begin{pmatrix}\cos(2\theta)&\sin(2\theta)\\
                  -\sin(2\theta)&\cos(2\theta)\end{pmatrix}.
\]
Its eigenvectors $(|g\rangle\pm i|b\rangle)/\sqrt2$ have eigenvalues
$e^{\pm2i\theta}$. Their squared coefficients in
$A|0\rangle=\sin\theta|g\rangle+\cos\theta|b\rangle$ are both $1/2$.

Set $M=2^m$ and prepare $m$ phase wires in their uniform superposition.
The binary controlled powers have exponents $1,2,\ldots,2^{m-1}$. They
implement $|j\rangle|v\rangle\mapsto|j\rangle Q^j|v\rangle$ using exactly
$\sum_{k=0}^{m-1}2^k=M-1$ controlled iterates. On an eigenvector of phase
$\omega\in[0,1)$, the phase state before the inverse Fourier transform is
\[
 |s_M(\omega)\rangle=M^{-1/2}\sum_{j=0}^{M-1}
 e^{2\pi i j\omega}|j\rangle.
\]
The inverse discrete Fourier transform gives amplitude
\[
 {1\over M}\sum_{j=0}^{M-1}e^{2\pi i j(\omega-y/M)}.
\]
For $d=\min_{k\in\mathbb Z}|\omega-y/M-k|>0$, its squared modulus is
\[
 D_M(y;\omega)={\sin^2(\pi M d)\over M^2\sin^2(\pi d)}.
\]
At an exactly representable phase the law is a unit atom at its phase
integer, interpreted by the continuous limit. Orthogonality of the two
logical eigenvectors removes cross terms when the logical register is
traced out. Thus the measured phase-register law is
\begin{equation}
 \Pr(Y=y)={D_M(y;\theta/\pi)+D_M(y;1-\theta/\pi)\over2}.
 \label{eq:qpe-dirichlet-mixture}
\end{equation}
The two nearest phase integers and this Dirichlet law yield the success
bound $8/\pi^2$ in \citet{brassard}, Theorem 11. With output
$\widetilde p=\sin^2(\pi Y/M)$, Theorem 12 gives
\begin{equation}
 |\widetilde p-p|\leq {2\pi\sqrt{p(1-p)}\over M}
                         +{\pi^2\over M^2}
 \quad\hbox{with probability at least }{8\over\pi^2}.
 \label{eq:qpe-published-amplitude-bound}
\end{equation}
Its error transfer can also be checked directly. For a phase-angle error
$h$ with $|h|\leq\pi/M$, periodicity and symmetry select an equivalent
angle representative, and
\[
 \sin^2(\theta+h)-\sin^2\theta
 =\sqrt{p(1-p)}\sin(2h)+(1-2p)\sin^2h.
\]
Using $|\sin(2h)|\leq2|h|$, $|\sin h|\leq|h|$ and
$|1-2p|\leq1$ gives the stated error. The boundary probabilities $p=0,1$
are treated separately by their exact eigenphases; for even $M$ their ideal
outputs are respectively zero and one. The numerical Dirichlet-law and
coherent-state sweeps are diagnostics of the implementation and do not
replace the published probability theorem.

\subsection{Finite plans and approximate inverse Fourier transform}

Let $\Pi=355/113$. The retained exact Machin-series enclosure independently
verifies $\pi<\Pi$. From the common proved prior define
\[
 v(u)=\max_{0\leq p\leq u}p(1-p)
 =\begin{cases}u(1-u),&u\leq1/2,\\1/4,&u>1/2.\end{cases}
\]
I compute a rational $s(u)\geq\sqrt{v(u)}$ by integer square-root and outward
rounding on a $2^{-100}$ lattice. Let $\xi=10^{-60}$ denote the reserved
classical output allowance. For an active scalar target $e\in[10^{-6},1/2)$,
I select the least supported power of two $M\geq2$ satisfying
\begin{equation}
 E_M(u)+\xi\leq e,
 \qquad E_M(u)={2\Pi s(u)\over M}+{\Pi^2\over M^2}.
 \label{eq:qpe-rational-binary-plan}
\end{equation}
Every comparison here is an exact rational comparison. If $u\leq2e$, the
whole-prior midpoint $u/2$ already has deterministic error at most $e$ and
the plan has no quantum calls. An active plan is never forced when this
zero-shot certificate applies. The declared code guard supports at most
32 phase wires; a wider plan would require a further Fourier compilation
certificate.

Let $U=F_M^{-1}$, let $W$ be the actually emitted Clifford+T
inverse Fourier tape and let $\varphi$ be its recorded aggregate scalar
correction. Define $\widehat U=e^{i\varphi}W$. The certified operator-norm
bound is $\|\widehat U-U\|\leq\delta\leq1/64$, with the specified register
ordering. Multiplication of the final uncontrolled tape by this unit scalar
has no effect on its measurement law and requires no physical phase gate. For any normalised joint input, the
two output pure states differ in vector norm by at most $\delta$; the
trace distance of their density operators is no larger. Partial trace and
measurement do not increase trace distance, so the probability of each
measurement event changes by at most $\delta$. Let $\eta_m$ be the actual retained norm upper bound for the chosen width.
For the confidence arithmetic I form the exact outward binary enclosure
$\delta_m=2^{-80}\lceil2^{80}\eta_m\rceil$. Since $1/64$ belongs to this
lattice, $\eta_m\leq\delta_m\leq1/64$. The difference is less than
$2^{-80}$; I retain $\eta_m$ separately in the Fourier profile. This
coarsening prevents unnecessarily large denominator powers in the finite
binomial record while keeping a rigorous success lower bound. Therefore
each emitted repetition succeeds with probability at least
\begin{equation}
 s_0(m)={8\over\Pi^2}-\delta_m
 \geq {8\over\Pi^2}-{1\over64}>{1\over2}.
 \label{eq:qpe-emitted-good-probability}
\end{equation}
The binary choice $M$ is determined before this confidence constant and does
not depend on $R$. Hence I may use the outward enclosure of its actual certificate in the
primary amplification without a circular synthesis or stopping choice. I
retain the $1/64$ substitution as a separate conservative envelope. In the
following formulae $s_0$ denotes the selected $s_0(m)$.
The norm allowance changes the failure probability; it is not charged
again as an amplitude bias. This transfer assumes exact logical controlled
iterates. Hardware noise, finite syndrome failure and gate faults require
additional budgets before a physical-device claim is made.

For phase outcome $y$, symmetry reduces its output to
$\sin^2(\pi\min(y,M-y)/M)$. The existing directed rational cosine
calculation returns an interval $[\ell_y,h_y]$ containing this value. Its
width is checked to be at most $2\xi$. The midpoint consequently differs
from the exact post-processing value by at most $\xi$. Projection of that
midpoint onto $[0,u]$ cannot increase its distance from the true $p$.
Combining this fact with \eqref{eq:qpe-rational-binary-plan} gives error at
most $e$ on every good emitted repetition.

\subsection{Fresh repetitions and exact confidence amplification}

I use a prescribed odd number $R$ of fresh complete repetitions. The
backend contract is that the repetitions are independent, each begins with
clean financial and phase registers and each has good-event probability
at least $s_0$. Reset and preparation are charged. A seeded ideal simulator
provides reproducible diagnostics but does not establish this independence
premise for a physical backend.

If at least $(R+1)/2$ estimates lie in $[p-e,p+e]$, their median lies in
the same interval. Independent bad-event indicators with probabilities at
most $1-s_0$ are stochastically dominated by independent Bernoulli variables
of parameter $1-s_0$: couple each indicator to its own independent uniform
variable and increase its threshold to $1-s_0$. The median-failure
probability is therefore at most the exact rational number
\begin{equation}
 B_R(s_0)=\sum_{j=(R+1)/2}^{R}{R\choose j}
       (1-s_0)^j s_0^{R-j}.
 \label{eq:qpe-exact-binomial-confidence}
\end{equation}
The primary finite plan chooses the smallest positive odd $R$ for which
$B_R(s_0)\leq\alpha$. Checking every smaller odd integer proves this finite minimality.
Termination follows from Hoeffding's inequality:
\[
 B_R(s_0)\leq\exp\{-2R(s_0-1/2)^2\}.
\]
The conservative comparison retains the least odd integer satisfying
$2R(s_0-1/2)^2\geq\overline{\log(1/\alpha)}$. This logarithm is also
enclosed with exact rational arithmetic. Write $x=2^k b$ with $1\leq b<2$;
for $z=(b-1)/(b+1)\leq1/3$,
\[
 \log b=2\sum_{j=0}^{N-1}{z^{2j+1}\over2j+1}+T_N,
 \qquad0\leq T_N\leq
 {2z^{2N+1}\over(2N+1)(1-z^2)}.
\]
The same formula encloses $\log2$ and hence $k\log2+\log b$. The retained
$N=80$ is a finite choice with the displayed positive remainder, not a
floating-point logarithm assumption. I retain both confidence rules, so the
stronger exact-binomial comparison does not inherit an unnecessary
Hoeffding repetition constant.

\subsection{Financial error allocation and complete resource charges}

For a range flag estimating $p=\mathbb E[X]/r$, a scalar probability radius
$e$ gives numerical mean radius $re$. For a clipped flag estimating
$p=\mathbb E[\min(X,r)]/r$, the checked tail enclosure
$0\leq\mathbb E[(X-r)^+]\leq D$ gives the reconstructed midpoint
$r\widehat p+D/2$ and mean radius $re+D/2$. The prescribed admissible choice
must satisfy
\[
 re+D/2\leq t,\qquad t=0.005/H.
\]
The original finite pricing certificate remains a separate deterministic
allowance at most $0.005$ in price currency units. Multiplication by $H$
therefore yields the retained undiscounted $0.01$ price radius on the
coverage event. Discounting is then applied exactly as in the existing
financial comparison. A clipping choice with $D/2\geq t$ is rejected unless
the common whole-prior argument already removes all estimation. For
dyadic buckets, the same scalar QPE guarantee can be used after the existing
weighted error allocation and failure-probability union bound; support,
moment and financial certificates remain common to competing methods.

Let $C_A$ denote any additive gate or internal-instrument metric of the actual emitted flag preparation,
$C_{\mathrm{CQ}}$ the corresponding metric of its emitted controlled
iterate, and $C_F(m)$ that of the certified inverse Fourier tape. The total before terminal phase readout and final reinitialisation
metric is
\begin{equation}
 R\{C_A+(M-1)C_{\mathrm{CQ}}+C_F(m)+mC_H\}.
 \label{eq:qpe-complete-charge}
\end{equation}
The final term is the phase-register preparation. The inverse Fourier
construction itself contains $m$ Hadamards, $m(m-1)/2$ controlled phases
and $\lfloor m/2\rfloor$ swaps before their explicit Clifford+T expansion.
Each swap is charged as three CNOT gates. Its controlled-phase tapes retain
the certified synthesis error and aggregate phase convention; these
rotations are not assigned zero T cost. There are $Rm$ final phase-register
observations and $R$ complete reinitialisations, in addition to all measured
AND observations, corrections and resets already in $C_A$ and
$C_{\mathrm{CQ}}$. Sequential repetitions and controlled iterates permit
scratch reuse. I add the $m$ simultaneously live phase wires and the
certified synthesis ancillas to the largest financial/reflection workspace.
Serial additive depth caps are conservative scheduling bounds; they do not
assert depth-optimal parallel execution.

The original unit-range, certified range, clipped and dyadic flags are
eligible only with their own complete profiles, identical financial
tolerance and common certified information. This canonical comparator
closes a particular missing baseline. It does not by itself establish
optimality against all amplitude estimators, priority of the finite
certificate, or a conversion from logical T counts to physical runtime.
